\documentclass[10pt]{article}

% Specify the margins and text size.
\setlength{\textwidth}{6.4in}
\setlength{\textheight}{9.5in}
\setlength{\oddsidemargin}{0pt}
\setlength{\evensidemargin}{0pt}
\setlength{\topmargin}{0pt}
\setlength{\hoffset}{.05in}
\setlength{\voffset}{-1in}

\setlength{\parskip}{5pt}
\setlength{\parindent}{0pt}

% Load some fonts and symbol packages
\usepackage{latexsym}
\usepackage{pifont}       % contains 'star' symbol for counterinsurgency handbook title
\usepackage{yfonts} 
\usepackage{amsmath}
\usepackage{amsfonts}

\usepackage{graphicx}     % actually, this is loaded by pstricks
\usepackage[T1]{fontenc}
\usepackage{ifthen}
\usepackage{pstricks,pst-grad,pst-text,pst-node,multido,pst-plot,calc,pst-3dplot}
%\usepackage[all]{xy}
%\usepackage{animate}

% The hyperref package inserts links into the pdf.
\definecolor{MyLinkColor}{rgb}{.1,.2,1}
\definecolor{MyCiteColor}{rgb}{.1,1,.2}
\definecolor{MyURLColor}{rgb}{.4,.4,.4}
\usepackage[backref=true,pagebackref=false,hyperindex,colorlinks=true,
  linkcolor=MyLinkColor,urlcolor=MyURLColor]{hyperref}


% The tweaklist package is something I found on the web.  It provides a simple interface
% for making changes to spacing used in the itemize and enumerate environments.  Comment
% this out if you don't care to use tweaklist.
\usepackage{tweaklist}
\renewcommand{\itemhook}{\setlength{\parskip}{2pt}\setlength{\parsep}%
{1pt}\setlength{\topsep}{0pt}\setlength{\itemsep}{0pt}}

\newcommand{\U}{\underline{\hspace{5pt}}}

\usepackage{listings}
\newcommand{\Z}{\hphantom{0}}

\begin{document}
\pagestyle{empty}
\lstset{language=R, showspaces=false, showstringspaces=false}

\href{http://www.su.edu}{\includegraphics[height=1.75cm]{sulogo.eps}}
\vspace{-1.69cm}

{{\ }\hfill\small
\begin{tabular}{cl}
& Math 207\\
& Introduction to Statistics\\
\end{tabular}
}
\setlength{\baselineskip}{1.05\baselineskip}
\bigskip

\begin{center}
\textbf{\large  Quiz 3}
\end{center}
\medskip

1. For the data set below, (a)~find the equation for the  regression line,
(b)~draw the scatter plot and the two line.\vspace{-15pt}
\begin{center}
\begin{tabular}{cc}
x & y\\ \hline
1 & 1\vphantom{\LARGE Y} \\
3 & 7\\
4 & 8\\
4 & 9\\
5 & 10\\
7 & 13
\end{tabular}
\end{center}

\vfill
\eject
{\ }\bigskip

2. In a study of the stability of IQ scores, a large group of individuals is
tested once at age 18 and again at age 35.  The following results are obtained.

\begin{center}
\begin{tabular}{llll}
age 18: & average score $\approx 100$, & $\mbox{sd}\approx 15$\\
age 35: & average score $\approx 100$, & $\mbox{sd}\approx 15$,  & $r\approx 0.8$\\
\end{tabular}
\end{center}

\hspace{20pt} a) Estimate the average score at age 35 for all the individuals who scored 115 at
 age 18.
\vspace{2in}

\hspace{20pt} b) Predict the score at age 35 for an individual who scored 115 at age 18.
\vspace{2in}

\hspace{20pt} c) Predict the score at age 35 for an individual who scored 130 at age 18.
\vfill
\eject

\end{document}

